\chapter{Module 2: Processing Pipeline}
\label{ch:module-processing}

\begin{tcolorbox}[
  enhanced,
  colback=LightGrey,
  colframe=MidGrey,
  fonttitle=\bfseries\sffamily,
  title={Module Information},
  sharp corners=south,
  boxrule=0.6pt
]
\begin{description}[font=\sffamily\bfseries, leftmargin=4cm, style=sameline]
  \item[Prerequisites] Module~1 (Chapter~\ref{ch:module-capture}); basic comfort with command-line interfaces
  \item[Duration] 4--6 hours (including processing wait time)
  \item[Hardware] Computer with NVIDIA GPU (8\,GB+ VRAM), or access to a cloud GPU instance
  \item[Software] Docker Desktop, web browser, terminal emulator
  \item[Budget tier] Free (with institutional GPU) or Standard (\pounds 1--3/hour cloud GPU)
\end{description}
\end{tcolorbox}

\medskip

\noindent\textbf{Learning Objectives.} Upon completing this module, participants will be able to:
\begin{enumerate}
  \item Install Docker and the \texttt{gaussian-toolkit} on a local workstation or cloud GPU instance.
  \item Upload captured video to the pipeline's web interface and initiate processing.
  \item Monitor pipeline progress through each stage (\acrshort{colmap}, training, optional mesh extraction).
  \item Interpret the key quality indicators in the pipeline output.
  \item Identify and resolve the five most common processing failures.
  \item Access and navigate the reconstructed scene in the built-in browser viewer.
\end{enumerate}

\section{Installing Docker and the Gaussian Toolkit}
\label{sec:proc-install}

The \texttt{gaussian-toolkit} is distributed as a set of Docker containers, as described in Chapter~\ref{ch:pipeline}. This eliminates the need to install CUDA, Python, \acrshort{colmap}, or any of the other GPU-accelerated dependencies manually---a process that, in our experience, takes even experienced developers 4--8 hours and is the single most common point of abandonment for new users of 3D reconstruction tools.

\subsection{Prerequisites}

Before beginning installation, verify that your system meets the following requirements:

\begin{itemize}
  \item \textbf{Operating system:} Linux (Ubuntu 20.04+), Windows 10/11 with WSL2, or macOS (CPU-only, significantly slower).
  \item \textbf{GPU:} NVIDIA GPU with 8\,GB+ VRAM and CUDA compute capability 7.0 or higher. The RTX 3060 (12\,GB) is the minimum recommended for practical use; the RTX 3080/4070 (16\,GB) or above is recommended for comfortable operation.
  \item \textbf{Disk space:} At least 50\,GB free for Docker images and working data.
  \item \textbf{RAM:} 16\,GB minimum, 32\,GB recommended.
  \item \textbf{Docker:} Docker Desktop (Windows/macOS) or Docker Engine (Linux), version 24.0 or later.
  \item \textbf{NVIDIA Container Toolkit:} Required on Linux; included in Docker Desktop for Windows with WSL2.
\end{itemize}

\subsection{Installation Steps}

\begin{lstlisting}[language=bash, caption={Installing the \texttt{gaussian-toolkit}.}, label={lst:install}]
# 1. Clone the repository
git clone https://github.com/dreamlab-ai/gaussian-toolkit.git
cd gaussian-toolkit

# 2. Pull the pre-built Docker images (approximately 30 GB)
docker compose pull

# 3. Start the services
docker compose up -d

# 4. Verify all services are running
docker compose ps

# 5. Open the web interface
# Navigate to http://localhost:8080 in your browser
\end{lstlisting}

\begin{technote}
The initial \texttt{docker compose pull} downloads approximately 30\,GB of container images. On a typical university network connection (100\,Mbps), this takes 30--45 minutes. Subsequent starts are near-instant as images are cached locally. For workshop delivery, pre-pulling images on all participant machines before the session begins is strongly recommended.
\end{technote}

\subsection{Cloud GPU Alternative}

For participants without local GPU access, cloud GPU instances provide an accessible alternative. The following providers have been tested with the \texttt{gaussian-toolkit}:

\begin{table}[htbp]
  \centering
  \caption{Cloud GPU options for \texttt{gaussian-toolkit} processing.}
  \label{tab:cloud-gpu}
  \begin{tabular}{@{}L{2.5cm}L{3cm}L{2.5cm}L{2cm}L{2.5cm}@{}}
    \toprule
    \textbf{Provider} & \textbf{GPU} & \textbf{VRAM} & \textbf{Cost/hr} & \textbf{Notes} \\
    \midrule
    RunPod       & RTX 4090     & 24\,GB & \$0.44 & Docker-native \\
    Vast.ai      & RTX 3090     & 24\,GB & \$0.25 & Variable pricing \\
    Lambda Labs  & A10G         & 24\,GB & \$0.60 & Persistent instances \\
    Google Colab & T4 (free) / A100 (Pro) & 16/40\,GB & \$0--12/mo & Notebook interface \\
    \bottomrule
  \end{tabular}
\end{table}

For cloud deployment, the \texttt{gaussian-toolkit} provides a one-command setup script:

\begin{lstlisting}[language=bash, caption={Cloud instance quick setup.}, label={lst:cloud-setup}]
# On a fresh cloud GPU instance (Ubuntu + NVIDIA drivers)
curl -fsSL https://get.docker.com | sh
sudo nvidia-ctk runtime configure --runtime=docker
sudo systemctl restart docker
git clone https://github.com/dreamlab-ai/gaussian-toolkit.git
cd gaussian-toolkit && docker compose up -d
\end{lstlisting}

\section{Uploading Video to the Web Interface}
\label{sec:proc-upload}

The \texttt{gaussian-toolkit} web interface provides a drag-and-drop upload mechanism accessible at \texttt{http://localhost:8080} (local) or the cloud instance's public IP address.

\subsection{Upload Process}

\begin{enumerate}
  \item Open the web interface in a browser (Chrome or Firefox recommended).
  \item Click \textbf{New Project} and enter a descriptive project name (e.g., ``North Gallery 2026-04-01'').
  \item Drag the captured video file onto the upload area, or click to browse. Supported formats: MP4 (H.264/H.265), MOV, AVI.
  \item Select the processing configuration:
    \begin{description}
      \item[Quick] 7,000 training iterations, \acrshort{colmap} exhaustive matching. Fastest results (20--40 minutes for a 3-minute video on RTX 4090), lower fidelity.
      \item[Standard] 30,000 iterations, \acrshort{colmap} sequential matching. Recommended for most captures (1--2 hours).
      \item[Archival] 50,000 iterations, \acrshort{colmap} exhaustive matching, all three training variants (\acrshort{3dgs}, \acrshort{mrnf}, \acrshort{mcmc}). Highest fidelity, longest processing time (3--6 hours).
    \end{description}
  \item Optionally enable mesh extraction (TSDF and/or MILo) if mesh outputs are required.
  \item Click \textbf{Start Processing}.
\end{enumerate}

\begin{recommendation}
For training purposes, use the \textbf{Quick} configuration for initial runs. The reduced iteration count produces a visibly lower-fidelity result, but it completes in 20--40 minutes rather than 1--2 hours, allowing participants to iterate on capture technique within a single workshop session. Once the capture technique is validated, reprocess the same footage with the \textbf{Standard} or \textbf{Archival} configuration for the final preservation artefact.
\end{recommendation}

\section{Monitoring the Pipeline}
\label{sec:proc-monitoring}

The web interface provides real-time progress monitoring through a multi-stage status display. Understanding what each stage does, and what constitutes normal versus abnormal behaviour, is essential for troubleshooting.

\subsection{Stage 1: Frame Extraction}

The pipeline extracts frames from the uploaded video using FFmpeg. The interface displays:

\begin{itemize}
  \item \textbf{Total frames extracted.} A 3-minute 4K video at 30\,fps produces approximately 5,400 frames. The pipeline extracts at a configurable rate (default: every second frame, yielding approximately 2,700 frames for a 3-minute video).
  \item \textbf{Frames rejected.} Blurred or near-duplicate frames are discarded by the Laplacian variance filter. A rejection rate above 20\% suggests the camera was moving too fast.
  \item \textbf{Final frame count.} Typically 1,500--3,000 frames for a 3--5 minute capture.
\end{itemize}

\subsection{Stage 2: COLMAP Structure-from-Motion}
\label{sec:proc-colmap}

This is the most time-consuming and failure-prone stage. \acrshort{colmap} processes the extracted frames to determine camera positions and produce a sparse 3D point cloud (see Section~\ref{sec:tech-sfm}).

The interface displays:

\begin{itemize}
  \item \textbf{Feature extraction progress.} Each frame is analysed for distinctive visual features (SIFT descriptors). This is parallelised across CPU cores and completes relatively quickly.
  \item \textbf{Feature matching progress.} Frames are compared pairwise to identify shared features. This is the slowest sub-stage, growing quadratically with frame count for exhaustive matching and linearly for sequential matching.
  \item \textbf{Registered images.} The number of frames successfully placed in 3D space. This is the single most important quality indicator.
  \item \textbf{Registration rate.} Registered images divided by total frames. The target is $\geq 90\%$.
  \item \textbf{Sparse points.} The number of 3D points in the reconstructed sparse cloud. Typical values range from 50,000 to 500,000 depending on scene complexity and capture thoroughness.
\end{itemize}

\begin{keyfinding}
The \acrshort{colmap} registration rate is the earliest and most reliable predictor of final reconstruction quality. A rate below 80\% almost always produces a reconstruction with visible holes or distortion. A rate below 50\% typically indicates a fundamental capture problem (motion blur, insufficient overlap) and the footage should be recaptured rather than reprocessed.
\end{keyfinding}

\subsection{Stage 3: Gaussian Splat Training}

Once \acrshort{colmap} completes, the Gaussian splat training begins. The interface displays:

\begin{itemize}
  \item \textbf{Current iteration} and progress bar (e.g., 15,000/30,000).
  \item \textbf{Training loss.} The photometric error between the rendered Gaussians and the training images. This should decrease rapidly in the first few thousand iterations, then gradually plateau.
  \item \textbf{Gaussian count.} The number of Gaussians in the scene, which increases through densification and decreases through pruning. Typical converged values range from 500,000 to 5,000,000 depending on scene complexity.
  \item \textbf{Preview renders.} Every 1,000 iterations, the interface renders a preview from a fixed viewpoint, allowing participants to watch the reconstruction emerge in real time.
\end{itemize}

\subsection{Stage 4: Mesh Extraction (Optional)}

If mesh extraction was enabled, one or both of the following processes run after training:

\begin{description}
  \item[TSDF fusion] Renders depth maps from the trained Gaussians at each camera position, then fuses them into a volumetric grid that is converted to a triangle mesh via marching cubes. Produces water-tight meshes but may lose fine detail.
  \item[MILo] A learned mesh extraction method that directly converts the Gaussian representation to a triangle mesh using a neural network. Produces higher-detail meshes but requires the MILo sidecar container.
\end{description}

\section{Understanding the Output}
\label{sec:proc-output}

Upon completion, the pipeline produces the following output files, accessible through the web interface's download panel or directly from the Docker volume:

\begin{table}[htbp]
  \centering
  \caption{Pipeline output files and their uses.}
  \label{tab:output-files}
  \begin{tabular}{@{}l l L{6cm}@{}}
    \toprule
    \textbf{File} & \textbf{Format} & \textbf{Use} \\
    \midrule
    \texttt{scene.ply}   & PLY (Gaussian) & The native Gaussian splat representation. Primary preservation artefact. Loadable in browser viewers and Gaussian-aware applications. \\
    \texttt{scene.splat}  & SPLAT binary   & Compressed format for browser viewers (antimatter15/splat, gsplat.js). Faster loading, same visual quality. \\
    \texttt{scene.glb}    & glTF Binary    & Triangle mesh (if mesh extraction enabled). Compatible with Blender, Unreal, Unity, and web-based 3D viewers (model-viewer). \\
    \texttt{scene.usda}   & USD ASCII      & Universal Scene Description (if enabled). The preferred interchange format for production 3D workflows and long-term archival. \\
    \texttt{cameras.json} & JSON           & Camera positions and intrinsics from \acrshort{colmap}. Useful for registering additional captures of the same space. \\
    \texttt{metadata.json} & JSON          & The complete capture and processing metadata record (see Module~5, Chapter~\ref{ch:module-metadata}). \\
    \bottomrule
  \end{tabular}
\end{table}

\begin{technote}
The \texttt{.ply} file containing the Gaussian splat representation is typically 50--200\,MB for a standard scene. The \texttt{.splat} file is a sorted, compressed version at approximately 60\% of the PLY size. The \texttt{.glb} mesh file varies widely---from 10\,MB for a simplified mesh to 500\,MB+ for a high-detail extraction---and contains texture maps baked from the Gaussian representation.
\end{technote}

\section{Viewing Results in the Browser}
\label{sec:proc-viewing}

The pipeline's built-in browser viewer is accessible immediately upon completion at \texttt{http://localhost:8080/viewer/\{project-id\}}. The viewer supports:

\begin{itemize}
  \item \textbf{Orbit navigation.} Click and drag to rotate around the scene centre. Scroll to zoom.
  \item \textbf{Fly-through navigation.} WASD keys for movement, mouse for look direction. This mode allows free exploration of the reconstructed space.
  \item \textbf{Quality comparison.} Toggle between the Gaussian splat and mesh representations (if both were generated) to observe the quality difference first-hand.
  \item \textbf{Gaussian count display.} Shows the number of Gaussians being rendered and the current frame rate. Useful for assessing whether the scene is within the viewer's comfortable rendering budget.
\end{itemize}

\section{Troubleshooting Common Issues}
\label{sec:proc-troubleshoot}

Table~\ref{tab:troubleshoot} documents the five most common processing issues encountered during the pilot project, their symptoms, likely causes, and remedies.

\begin{table}[htbp]
  \centering
  \caption{Common processing issues and remedies.}
  \label{tab:troubleshoot}
  \begin{tabular}{@{}L{2.5cm}L{3cm}L{3cm}L{4cm}@{}}
    \toprule
    \textbf{Symptom} & \textbf{Likely Cause} & \textbf{Evidence} & \textbf{Remedy} \\
    \midrule
    \acrshort{colmap} registration $< 50\%$
      & Motion blur or insufficient frame overlap
      & Blurred frames visible when reviewing extracted images
      & Recapture at slower pace; increase frame extraction rate \\
    \addlinespace
    Large holes in reconstruction
      & Area not captured from enough viewpoints
      & Missing coverage visible in \acrshort{colmap} sparse point cloud
      & Capture additional footage of the missing area and merge \\
    \addlinespace
    Floating blobs (``floaters'')
      & Moving objects in scene or strong reflections
      & Floaters correspond to locations of people, screens, or mirrors
      & Recapture without moving elements; use SAM3 cleanup (Module~4) \\
    \addlinespace
    Out of GPU memory
      & Too many frames or excessively high resolution
      & CUDA out-of-memory error in training logs
      & Reduce frame extraction rate; downsample to 2K; use GPU with more VRAM \\
    \addlinespace
    Pipeline hangs at ``Feature matching''
      & Exhaustive matching on too many frames
      & Matching progress stalls at high frame counts ($>$3,000)
      & Switch to sequential matching mode; reduce frame count \\
    \bottomrule
  \end{tabular}
\end{table}

\section{Practical Exercise: Process Your Captured Video}
\label{sec:proc-exercise}

\begin{tcolorbox}[
  enhanced,
  colback=SuccessGreen!8,
  colframe=SuccessGreen,
  fonttitle=\bfseries\sffamily,
  title={Exercise 2.1: Pipeline Processing},
  sharp corners=south,
  boxrule=0.6pt
]
\textbf{Objective:} Process the video captured in Exercise~1.1 (Section~\ref{sec:cap-exercise}) through the full \texttt{gaussian-toolkit} pipeline.

\textbf{Time:} 30 minutes setup + 20--40 minutes processing (Quick mode) or 1--2 hours (Standard mode)

\textbf{Equipment:} Computer with Docker and \texttt{gaussian-toolkit} installed (per Section~\ref{sec:proc-install}), or cloud GPU instance

\textbf{Steps:}
\begin{enumerate}
  \item Upload your captured video to the web interface and select \textbf{Quick} configuration.
  \item While processing runs, record the following metrics as they appear:
    \begin{itemize}
      \item Total frames extracted: \rule{3cm}{0.4pt}
      \item Frames rejected (blur): \rule{3cm}{0.4pt}
      \item \acrshort{colmap} registration rate: \rule{3cm}{0.4pt}
      \item Sparse points: \rule{3cm}{0.4pt}
      \item Final Gaussian count: \rule{3cm}{0.4pt}
    \end{itemize}
  \item When processing completes, open the browser viewer and navigate through your reconstruction.
  \item Identify and note any quality issues: holes, floaters, distortion, missing areas.
  \item Correlate any quality issues with your capture technique: can you identify what you would do differently?
  \item If your registration rate was below 80\%, recapture and reprocess.
\end{enumerate}

\textbf{Expected outcome:} A viewable Gaussian splat reconstruction of the room captured in Exercise~1.1, accompanied by an understanding of how capture decisions map to reconstruction quality. Participants should be able to articulate at least one specific improvement they would make to their capture technique.
\end{tcolorbox}

\begin{tcolorbox}[
  enhanced,
  colback=SuccessGreen!8,
  colframe=SuccessGreen,
  fonttitle=\bfseries\sffamily,
  title={Exercise 2.2: Quality Comparison},
  sharp corners=south,
  boxrule=0.6pt
]
\textbf{Objective:} Compare Quick and Standard processing of the same footage to understand the quality--time trade-off.

\textbf{Time:} 1--2 hours (processing time for Standard configuration)

\textbf{Steps:}
\begin{enumerate}
  \item Reprocess the same video using the \textbf{Standard} configuration (30,000 iterations).
  \item Open both results in the browser viewer (two tabs) and compare:
    \begin{itemize}
      \item Fine detail (text on labels, surface textures)
      \item Edge sharpness (transitions between objects)
      \item Floaters and artefacts
      \item Overall colour fidelity
    \end{itemize}
  \item Record your observations: is the quality difference worth the additional processing time for this particular scene?
\end{enumerate}

\textbf{Expected outcome:} Participants should observe noticeably finer detail and fewer artefacts in the Standard output, and be able to make an informed decision about which configuration to use for different capture scenarios.
\end{tcolorbox}
